\documentclass[11pt]{article}
\usepackage[utf8]{inputenc}
\usepackage{amsmath,amssymb}
\usepackage[margin=1in]{geometry}
\usepackage{hyperref}
\emergencystretch=3em

\title{The coefficient moments as absolutely convergent series}
\author{Claude, with Daniel Murfet}
\date{2026-09-07}

\begin{document}
\maketitle
\sloppy

\begin{abstract}
Completion of Astra~\#8 rank~4. Two interchanges make the $d=2$ Taylor-tree coefficients
explicit. (1) The face moments: $M[\alpha;\ell;U_{\alpha_i}]=k_1^{-1}\sum_jq_{\gamma_i}(j)\,M[\alpha;\ell;c_{ij}]$,
absolutely convergent, by the weighted $L^1$ interchange under the envelope
$|c_{ij}(s)|\le H(s)\rho^{-i-j}$, $H(s)\le C_0(1+s)^De^{\beta sL}$, with the weights
$\sum_j|q_\gamma(j)|\rho^{-j}$ summable. (2) The Gaussian fluctuation moments: for the amplitude
array, $c_k(s)=e^{\beta sx_{00}}\sum_{n\le|k|}\frac{(\beta s)^n}{n!}(y*a^{*n})_k$ is a FINITE sum (degree
support), hence
$M[\alpha;\ell;c_k]=\sum_{n\le|k|}\frac{\beta^n}{n!}(y*a^{*n})_k\int_0^\infty s^{\alpha+n-1}(\log s)^\ell e^{-\beta s^2+\beta sx_{00}}ds$,
the paper's moments $\int t^{p/2}e^{-\beta t+\beta\sqrt t\,\xi(0)}dt$ in the variable $s=\sqrt t$.
Zero \texttt{sorry}/\texttt{axiom}.
\end{abstract}

\section{The things formalised}
{\small
\begin{verbatim}
theorem axisPrim_div_summable (γ b ρ : ℝ) (hb : 0 < b) (hbρ : b < ρ) :
    Summable fun j : ℕ => |axisPrim γ b j| / ρ ^ j
theorem logMoment_canonU_series (β b ρ C₀ L : ℝ) (h₁ h₂ k₁ k₂ D : ℕ) (hβ : 0 < β)
    (hb : 0 < b) (hbρ : b < ρ) (hk₁ : 0 < k₁) (c : ℕ × ℕ → ℝ → ℝ)
    (hcc : ∀ ij, Continuous (c ij)) (H : ℝ → ℝ)
    (hc : ∀ ij s, |c ij s| * ρ ^ (ij.1 + ij.2) ≤ H s) (hC₀ : 0 ≤ C₀)
    (henv : ∀ s, 0 ≤ s → H s ≤ C₀ * (1 + s) ^ D * Real.exp (β * s * L)) (ℓ : ℕ) (i : ℕ) :
    (Summable fun j : ℕ => axisPrim ((k₂ : ℝ) * uExp h₁ k₁ i - h₂ - 1) b j
        * logMoment β (uExp h₁ k₁ i) ℓ (c (i, j))) ∧
      logMoment β (uExp h₁ k₁ i) ℓ
          (canonU b h₁ h₂ k₁ k₂ (anaFaceU c b) (fun i j s => c (i, j) s) (uExp h₁ k₁ i))
        = 1 / (k₁ : ℝ) * ∑' j : ℕ, axisPrim ((k₂ : ℝ) * uExp h₁ k₁ i - h₂ - 1) b j
            * logMoment β (uExp h₁ k₁ i) ℓ (c (i, j))
theorem ampCoeff_eq_sum (β : ℝ) (x y : ℕ × ℕ → ℝ) (k : ℕ × ℕ) (s : ℝ) :
    ampCoeff β x y k s = Real.exp (β * s * x (0, 0))
      * ∑ n ∈ Finset.range (k.1 + k.2 + 1),
          (β * s) ^ n / (n.factorial : ℝ) * conv y (convPow (dropConst x) n) k
theorem gaussMoment_integrableOn (β α x₀ : ℝ) (hβ : 0 < β) (hα : 0 < α) (ℓ n : ℕ) :
    IntegrableOn (fun s => s ^ (α - 1) * Real.log s ^ ℓ
      * (Real.exp (-β * s ^ 2) * (s ^ n * Real.exp (β * s * x₀)))) (Ioi 0)
theorem logMoment_ampCoeff_eq_sum (β α : ℝ) (hβ : 0 < β) (hα : 0 < α) (x y : ℕ × ℕ → ℝ)
    (k : ℕ × ℕ) (ℓ : ℕ) :
    logMoment β α ℓ (ampCoeff β x y k)
      = ∑ n ∈ Finset.range (k.1 + k.2 + 1),
          β ^ n / (n.factorial : ℝ) * conv y (convPow (dropConst x) n) k
            * logMoment β α ℓ (fun s => s ^ n * Real.exp (β * s * x (0, 0)))
\end{verbatim}
}

\section{Mathematics}

For the face moments the term family is $f_j(s)=c_{ij}(s)q_{\gamma_i}(j)$ with
$|f_j(s)|\le\frac{H(s)}{\rho^i}\cdot\frac{|q_{\gamma_i}(j)|}{\rho^j}$; the weights are summable because
$|q_\gamma(j)|\le b^jb^{-\gamma}/(J-\gamma)$ for $j\ge J>\gamma$ gives a geometric tail
$(b/\rho)^j$. So $\sum_j|f_j(s)|\le\frac{C_0Q}{\rho^i}(1+s)^De^{\beta sL}$ with $Q=\sum_j|q(j)|\rho^{-j}$,
an envelope of the shape accepted by \texttt{hasSum\_logMoment\_of\_envelope} (unit~98), and the
interchange follows. For the Gaussian moments, $E_m(t)=\sum_{n<M}t^n/n!\,(a^{*n})_m$ for any
$M>|m|$ (the extra terms vanish by degree support), so the convolution $(y*E(t))_k$ can be summed
over a common range $n\le|k|$ and the two finite sums commute; the finite sum then passes through
the $s$-integral once each term $s^{\alpha-1}(\log s)^\ell e^{-\beta s^2}s^ne^{\beta sx_{00}}$ is
integrable, which is the envelope integrability of unit~96 with $R=n$, $a'=x_{00}$.

\section{Paper-facing meaning}

Combining units 121--126, every coefficient of the $d=2$ Taylor tree (equal starts) is now an
explicit absolutely convergent series in the convolution data $(y*a^{*n})_{ij}$ of the coefficient
arrays of $\eta$ and $\xi-\xi(0)$, weighted by $b^{j-\gamma}/(j-\gamma)$ (or $\log b$) and by the
Gaussian moments $\int_0^\infty s^{\alpha+n-1}(\log s)^\ell e^{-\beta s^2+\beta s\xi(0)}ds$. This is the
content of the paper's ``coefficients are absolutely convergent series expressed as derivatives of
$\xi$, $\eta$ and the fluctuation function'' once the coefficient arrays are identified with Taylor
data (the 1D identification exists; the 2D one is Taylor-coefficient uniqueness for double series).
What remains for the paper's $d=2$ theorem is the unequal-starting-exponent case (Astra~\#8
rank~3), which needs a new corner formula and an audit of the face expansion for negative $\gamma$.

\section{Lean notes}

\begin{itemize}
\item \texttt{HasSum.congr\_fun} wants the equation in the direction ``new term $=$ old term''.
\item \texttt{Finset.range\_subset} is no longer an iff usable with \texttt{.2}; use
\texttt{Finset.range\_mono}.
\item \texttt{Finset.sum\_comm} needs the inner sum exposed: push the outer factor in with
\texttt{Finset.mul\_sum} inside the \texttt{sum\_congr} first.
\item \texttt{gcongr} may leave a single arithmetic side goal (here $s\le1+s$) rather than one per
factor; bullets then fail with ``no goals''.
\item \texttt{positivity} cannot see $0\le s^{\alpha-1}$ from $s\in\mathrm{Ioi}\,0$; pass
\texttt{Real.rpow\_nonneg hs.le \_}.
\item \texttt{integral\_finset\_sum} is deprecated for \texttt{integral\_finsetSum}.
\item Writing \texttt{div\_eq\_mul\_inv b \^{} j} parses as \texttt{(div\_eq\_mul\_inv b) \^{} j};
prove the rpow identity as a separate \texttt{have}.
\end{itemize}

\end{document}
